\chapter{Human Parainfluenza Viruses (HPIVs)}

\section*{Definition and Overview}
Human parainfluenza viruses (HPIVs) are a group of four common respiratory viruses that cause illness across all ages. They are best known for causing croup in young children, a barking cough with breathing difficulty, and they also cause colds, bronchiolitis, and pneumonia. HPIVs are spread by droplets and contact, circulate mainly in autumn and winter, and reinfection occurs throughout life because immunity is short-lived. Most infections are mild, but infants, older adults, and immunocompromised people can develop serious lower respiratory disease.

\section*{Epidemiology}
HPIVs infect almost all children by age five, and they are a leading cause of respiratory hospitalisation in young children. Type 1 is the classic cause of croup, and types 1--3 cause seasonal epidemics in autumn and winter. HPIVs account for a significant share of childhood colds, croup, and pneumonia, and for respiratory illness and pneumonia in older adults. There is no vaccine, and outbreaks occur in childcare, schools, and healthcare facilities.

\section*{Aetiology and Pathophysiology}
\begin{itemize}
    \item \textbf{Types:} HPIV types 1, 2, 3, and 4, each with different age patterns and seasonality
    \item \textbf{Transmission:} respiratory droplets and contact with contaminated surfaces
    \item \textbf{Pathophysiology:} the virus infects the lining of the upper and lower airways, causing inflammation and swelling
    \item \textbf{Croup:} swelling of the upper airway (larynx and trachea) narrows the airway, causing a barking cough and noisy breathing
    \item \textbf{Bronchiolitis and pneumonia:} infection of the small airways and lung tissue in infants and at-risk adults
    \item \textbf{Immunity:} short-lived, so reinfection is common throughout life
\end{itemize}

\section*{Clinical Features}
\begin{itemize}
    \item Runny nose, cough, sore throat, and mild fever
    \item Croup in children aged 6 months to 3 years: barking cough, hoarse voice, and noisy breathing on inspiration (stridor), often worse at night
    \item Bronchiolitis in infants: cough, wheeze, and rapid breathing
    \item Pneumonia: fever, cough, and breathlessness
    \item Older adults: worsening of COPD and asthma, and pneumonia
    \item Ear infections and sinusitis
    \item Symptoms typically last 3--7 days
\end{itemize}

\section*{Differential Diagnosis}
\begin{itemize}
    \item Other respiratory viruses: RSV, influenza, rhinovirus, adenovirus, HMPV, and SARS-CoV-2
    \item Croup is distinguished from epiglottitis, a rare bacterial emergency with drooling and severe distress
    \item Bacterial pneumonia and whooping cough
    \item Asthma and allergies
    \item Testing identifies HPIV, though treatment is similar for most viral respiratory infections
\end{itemize}

\section*{When to Seek Medical Care}
Most HPIV infections are managed at home. See a doctor for breathing difficulty, a croupy cough that is severe or not settling, dehydration, or fever beyond a few days, and for infants and at-risk adults who become unwell.

\textbf{Call 000 (or local emergency services) for:}
\begin{itemize}
    \item Severe breathing difficulty, stridor at rest, or blue lips
    \item A child who cannot speak or is very agitated from breathing effort
    \item Drooling, inability to swallow, or a toxic appearance (possible epiglottitis)
    \item A baby who is floppy, lethargic, or not feeding
    \item High fever with confusion or drowsiness
    \item Severe dehydration
\end{itemize}

\section*{Diagnosis and Investigations}
\begin{itemize}
    \item \textbf{Clinical diagnosis:} the barking cough of croup and the winter respiratory pattern are usually sufficient
    \item \textbf{PCR swabs:} nose and throat swabs confirm the virus when needed
    \item \textbf{Chest X-ray:} for suspected pneumonia
    \item \textbf{Oxygen monitoring:} in hospitalised patients
    \item \textbf{Laryngoscopy:} rarely needed to exclude other causes of stridor
\end{itemize}

\section*{Management}
\begin{itemize}
    \item \textbf{Supportive care:} rest, fluids, and paracetamol for fever
    \item \textbf{Croup:} calm the child, sit them upright, and use cool mist or cool air; severe croup is treated with steroids and sometimes nebulised adrenaline
    \item \textbf{Bronchiolitis and pneumonia:} oxygen and breathing support in hospital if needed
    \item \textbf{No specific antiviral:} treatment is supportive, with antibiotics only for bacterial complications
    \item \textbf{Chronic disease care:} adjust asthma and COPD treatment during infection
    \item \textbf{Infection control:} hand hygiene, staying home while unwell, and covering coughs
\end{itemize}

\section*{Complications}
\begin{itemize}
    \item Severe croup requiring hospital care and occasionally intubation
    \item Bronchiolitis and pneumonia, especially in infants
    \item Respiratory failure in high-risk patients
    \item Exacerbations of asthma and COPD
    \item Ear infections and sinusitis
    \item Secondary bacterial pneumonia
    \item Outbreaks in childcare, aged care, and hospitals
\end{itemize}

\section*{Prognosis and Outcomes}
Most HPIV infections are mild and self-limiting, and croup resolves with supportive care and, when needed, steroid treatment. Infants and older adults with lower respiratory disease usually recover with hospital support. Reinfection is common but usually milder. There is no vaccine yet, so hygiene and protection of high-risk groups remain the main preventive measures.

\section*{Prevention and Self-Care}
\begin{itemize}
    \item Practise good hand hygiene, especially in winter
    \item Cover coughs and sneezes with the elbow
    \item Stay home when unwell
    \item Clean frequently touched surfaces
    \item Comfort a child with croup: stay calm, sit them upright, and use cool air
    \item Keep children's immunisations, including against flu, up to date
    \item Seek early medical care for infants and at-risk adults
    \item Follow public health advice during outbreaks
\end{itemize}

\section*{Sources and Further Reading}
The following reputable sources provide further information for healthcare students and patients seeking reliable, current advice about human parainfluenza viruses.

\begin{itemize}
    \item Healthdirect. \textit{Parainfluenza viruses.} https://www.healthdirect.gov.au/
    \item Royal Children's Hospital Melbourne. \textit{Croup fact sheet.} https://www.rch.org.au/
    \item NSW Health. \textit{Parainfluenza virus fact sheet.} https://www.health.nsw.gov.au/
    \item Centers for Disease Control and Prevention. \textit{Human parainfluenza viruses.} https://www.cdc.gov/
    \item NHS. \textit{Croup.} https://www.nhs.uk/conditions/croup/
    \item World Health Organization. \textit{Respiratory viruses information.} https://www.who.int/
\end{itemize}
